\chapter{Users, privileges, transport security, and encryption}
\label{ops:security}

\section{Understand the separate security boundaries}

MiniSQL has several complementary controls. Database authentication identifies
a user. Authorization decides which SQL operations that user can perform.
Encrypted transport protects traffic between a client and listener. Encryption
at rest protects selected persisted database artifacts when copied without
their external key. Operating-system permissions protect the files, keys,
process, and administrative launchers. Enabling one of these controls does
not implicitly enable the others.

A trusted-local server intentionally bypasses database authorization. It
binds to loopback, but another process running locally may still connect.
Use that mode only when local access is the intended administrative boundary.
For remotely reachable deployments, select native TLS, authenticate clients,
grant the minimum useful privileges, and protect the server's host and keys.

\begin{longtable}{p{.29\linewidth}p{.64\linewidth}}
\toprule
Control & What it does, and what remains outside its scope \\
\midrule
\endhead
User/password authentication & Establishes the database principal; this implementation then activates MiniSQL protected framing. Native TLS additionally provides certificate-based server identity. \\
Roles and grants & Restrict supported database/table operations. They do not
stop an operating-system account from reading database files it can access. \\
Native TLS & Authenticates the transport endpoint according to the selected
trust policy and encrypts traffic. It does not encrypt stored log files. \\
MiniSQL TDE & Encrypts supported paged files, WAL records, and temporary spill
records. It does not cover all filesystem metadata, text logs, or memory. \\
Encrypted backup & Authenticates and encrypts captured backup files with a
separate key. It cannot help if the key is lost or stored alongside the stolen
backup. \\
OS ACLs and volume encryption & Protect the installation and additional
artifacts. They do not replace application-level privilege design. \\
\bottomrule
\end{longtable}

\section{Bootstrap the administrator}
\index{administrator}\index{passwords}

A newly created database has built-in \code{admin} and \code{public}
principals. The administrator is an enabled built-in superuser. The public
role is not a login account; its membership is implicit for ordinary users.
The administrator must have a usable password before password-authenticated
connections are useful.

Stop the server and set the initial password locally:

\begin{lstlisting}
minisqld.exe --set-admin-password C:\MiniSQL-Data\db_ACTUAL_ID
\end{lstlisting}

The program prompts twice with echo disabled. On Linux use
\path{minisqld --set-admin-password /srv/minisql/db_ACTUAL_ID}.
The same local mechanism resets an existing user's password:

\begin{lstlisting}
minisqld.exe --set-user-password C:\MiniSQL-Data\db_ACTUAL_ID analyst
\end{lstlisting}

These commands open the database directly and require its storage lock.
Stop the server before using them. They are a local administrative recovery
mechanism, not a network password-reset endpoint. The caller also needs access
to the current encryption key when TDE is enabled.

Passwords contain 12 to 1024 UTF-8 bytes. Count bytes rather than assuming
that every visible character is one byte. New credentials use a 16-byte salt
and PBKDF2-HMAC-SHA-256 with 600,000 iterations, then persist the SCRAM-style
StoredKey/ServerKey pair rather than the salted password-equivalent secret.
The challenge/response transcript uses fresh nonces. Unknown and disabled
accounts use fake database-bound material and share the authentication-failure
response. Historical verifier records are upgraded when their password is
changed.

These implementation properties do not make the password exchange a PAKE.
Use TLS for remote traffic. Prompt-based commands avoid placing the password
in the operating system's process argument list, and mutable client secret
buffers are wiped after use. A password typed as a SQL literal can still be
recorded by an enabled SQL binlog or client history.

\sourcecode{src/apps/minisqld/main.ml}{setUserPassword and prompt-based local administration.}
\sourcecode{src/minisql/common/uuid.ml}{Password derivation, verifier material, and secret handling.}
\sourcecode{src/minisql/catalog/catalog.ml}{Built-in principals and stored credentials.}

\section{Create users and reusable roles}
\index{users}\index{roles}\index{CREATE USER}\index{CREATE ROLE}

The following demonstration uses a deliberately visible example password.
Replace it with a unique secret in a real installation, and decide beforehand
whether SQL history or statement logging is enabled.

\begin{lstlisting}[language=SQL]
CREATE ROLE reporting;
CREATE USER analyst WITH PASSWORD 'Replace-this-example-2026!';
GRANT CONNECT ON DATABASE TO analyst;
GRANT SELECT ON TABLE customers TO reporting;
GRANT SELECT ON TABLE orders TO reporting;
GRANT SELECT ON TABLE order_items TO reporting;
GRANT reporting TO analyst;
\end{lstlisting}

The role groups permissions by job. A second reporting user can receive the
same role without duplicating every table grant. Creating a user does not
automatically grant \code{CONNECT}. A correct password without the required
connection privilege is not enough. Likewise, table SELECT grants do not
implicitly grant database creation or maintenance.

Connect as the new user and test both an allowed SELECT and a deliberately
disallowed change on a disposable row. Testing the negative case helps detect
an accidentally broad grant, especially one inherited through \code{public}
or another role.

\subsection{Complete principal-management forms}

\begin{lstlisting}[language=SQL]
CREATE USER user_name WITH PASSWORD 'example-password';
CREATE ROLE role_name;
ALTER USER user_name WITH PASSWORD 'new_password';
ALTER USER user_name ENABLE;
ALTER USER user_name DISABLE;
DROP USER [IF EXISTS] user_name;
DROP ROLE [IF EXISTS] role_name;
\end{lstlisting}

Creating and dropping principals, and enabling or disabling accounts, require
database administration. A user may change its own password. The built-in
administrator cannot be disabled, and built-in principals are protected from
ordinary removal. Role-membership cycles are rejected. Public membership is
implicit and must not be assigned as an ordinary explicit role edge.

All these statements are autocommit-only. Do not put account setup in a
\code{BEGIN}/\code{COMMIT} block expecting transactional rollback of the
security catalog.

\sourcecode{src/minisql/sql/parser.ml}{parseCreatePrincipal, parseAlter, and principal DROP forms.}
\sourcecode{src/minisql/executor/executor.ml}{executeCreatePrincipal, executeAlterUser, and executeDropPrincipal.}

\section{The exact privilege model}
\index{privileges}\index{GRANT}

\begin{longtable}{p{.25\linewidth}p{.68\linewidth}}
\toprule
Scope / privilege & Meaning in the implementation \\
\midrule
\endhead
Database: \code{CONNECT} & Allows an authenticated user to establish access
to the database. \\
Database: \code{CREATE} & Permits table creation and supported schema/view/
sequence definition operations. Procedure creation requires CREATE ROUTINE. \\
Database: \code{MAINTAIN} & Permits ANALYZE, VACUUM, and REINDEX. \\
Database: \code{ADMIN} & Permits security administration, operational status,
process inspection, administrative cancellation, and shutdown. \\
Database: \code{CREATE ROUTINE} & Creates stored procedures. \\
Database: \code{SHOW\_ROUTINE} & Reveals routine source definitions as well as metadata. \\
Database: \code{SET\_ANY\_DEFINER} & Permits selecting another existing ordinary
user as definer; choosing an administrator still requires ADMIN. \\
Database: \code{SELECT} & Reads database tables and reveals routine source
definitions, subject to the ordinary supported object rules. \\
Database/procedure: \code{EXECUTE} & Calls every procedure in the database,
or one named procedure. \\
Database/procedure: \code{ALTER ROUTINE} & Alters, replaces, or drops every
procedure in the database, or one named procedure. \\
Table: \code{SELECT} & Reads table data; required for DESCRIBE and
SHOW INDEXES on that table. \\
Table: \code{INSERT} & Inserts data. INSERT SELECT also requires read access
to its source tables. \\
Table: \code{UPDATE} & Updates table data. \\
Table: \code{DELETE} & Deletes data and authorizes TRUNCATE. \\
Table: \code{REFERENCES} & A recognized table privilege in the security
catalog. Do not infer complete SQL-standard referential authorization from
its presence; validate the concrete schema workflow. \\
Table: \code{INDEX} & Creates or drops explicit indexes on the table. \\
Table: \code{ALTER} & Alters the table and permits trigger creation on it,
subject to body authorization. \\
Table: \code{DROP} & Drops the table. \\
\bottomrule
\end{longtable}

\begin{lstlisting}[language=SQL]
GRANT privilege [, privilege ...] ON DATABASE
  TO principal [WITH GRANT OPTION];
GRANT privilege [, privilege ...] ON [TABLE] table_name
  TO principal [WITH GRANT OPTION];
GRANT EXECUTE ON PROCEDURE procedure_name
  TO principal [WITH GRANT OPTION];
GRANT ALTER ROUTINE ON PROCEDURE procedure_name TO principal;
GRANT CREATE ROUTINE ON DATABASE TO principal;
GRANT SHOW_ROUTINE ON DATABASE TO principal;
GRANT SET_ANY_DEFINER ON DATABASE TO principal;
GRANT ALL [PRIVILEGES] ON DATABASE TO principal;
GRANT ALL [PRIVILEGES] ON TABLE table_name TO principal;
GRANT role_name TO principal [WITH ADMIN OPTION];
\end{lstlisting}

\code{ON DATABASE} means the currently opened database and does not take a
database name. \code{ALL} expands to the implemented privileges for the
selected scope. Only privileges implemented for that scope are accepted;
SELECT, EXECUTE and ALTER ROUTINE have the documented database-wide forms.
The parser
accepts one grantee per statement, so grant to multiple users with separate
statements.

\code{WITH GRANT OPTION} permits delegation of an object privilege.
\code{WITH ADMIN OPTION} permits delegation of a role membership. They
are not interchangeable. Database administrators can manage these grants;
other users need the appropriate delegation right. A role can contain other
roles as long as the membership graph remains acyclic.

Grants to \code{public} affect every user through implicit membership. A
convenient public read grant can therefore defeat an intended per-user
restriction. Review this role alongside direct grants and nested memberships.

\subsection{Views, procedures, triggers, and metadata}

Views are authorized through their underlying query sources. They are not
documented as a security-definer mechanism that lets a caller bypass base-table
permissions. Creating a view requires \code{CREATE} and authorization of its
query. Calling a procedure requires \code{EXECUTE} on that procedure or on the
current database. New MySQL-style definitions default to \code{SQL SECURITY
DEFINER}: statements in the body run with the persisted definer identity.
\code{SQL SECURITY INVOKER} uses the caller instead. Earlier AS definitions
retain an INVOKER default. The definer must also hold EXECUTE when its routine
is called; merely owning the definition does not bypass that check.
SET\_ANY\_DEFINER permits selecting another existing ordinary account.
Selecting an administrator as definer requires administrative permission.

With the default \code{automatic\_sp\_privileges = 1}, creation grants the
creator EXECUTE and ALTER ROUTINE without GRANT OPTION. An administrator may
set \code{SET SESSION automatic\_sp\_privileges = 0} to suppress these automatic
grants for later creations in that connection. This is MiniSQL's session
setting, not MySQL's server-global configuration mechanism. If automatic
grants are disabled, grant the required EXECUTE privilege explicitly.

A procedure's durable identifier anchors its grants. Dropping and recreating
the same name allocates a fresh identifier, so old grants cannot authorize the
replacement even if interrupted cleanup leaves a stale grant record. Definer
identity also uses a persisted principal identifier; deleting and recreating
a same-named account does not adopt its predecessor's routines. Calls fail if
the required definer account is missing or disabled.

Creating a trigger requires \code{ALTER}
on its target and authorization of its body. Dropping or altering a trigger
uses database \code{CREATE} in the present dispatcher.

Routine existence and parameter metadata are visible to administrators,
the definer, and principals with the relevant routine or source-reading
privileges. Reading the definition itself requires the definer identity,
ADMIN, SHOW\_ROUTINE or database SELECT. EXECUTE, CREATE ROUTINE and ALTER
ROUTINE alone do not reveal its source: SHOW CREATE and ROUTINES return a
NULL source field where the metadata is visible but source access is absent.
Other
\code{SHOW TABLES} and virtual \code{information\_schema} queries are not
a general promise of privilege-filtered object-name secrecy. In contrast,
\code{DESCRIBE} and \code{SHOW INDEXES} require SELECT on their target
table. Design authorization around supported operations, and test metadata
visibility if hiding object names is a requirement.

Procedure CREATE, ALTER and DROP validate authorization before their implicit
transaction commit. A denied definition change therefore cannot commit a
caller's pending DML. An authorized operation may still fail later semantic
validation after that commit; follow the transaction chapter's boundary rules.

\sourcecode{src/minisql/executor/executor.ml}{privilegeCode, allPrivilegeCodes, authorizeStatement, authorizeDefinitionStatement, and authorizeNamedSource.}

\section{Revoke privileges and delegated memberships}
\index{REVOKE}

\begin{lstlisting}[language=SQL]
REVOKE privilege [, privilege ...] ON DATABASE
  FROM principal [RESTRICT | CASCADE];
REVOKE privilege [, privilege ...] ON [TABLE] table_name
  FROM principal [RESTRICT | CASCADE];
REVOKE role_name FROM principal [RESTRICT | CASCADE];
\end{lstlisting}

The default behavior is restrictive. If another grant depends on the
delegated right, a restrictive revoke reports the dependency instead of
silently deleting a wider chain. \code{CASCADE} removes affected dependent
grants as well. There is no separate \code{REVOKE GRANT OPTION FOR}
grammar in the current parser.

For example, if Alice received SELECT with grant option and granted it to
Bob, revoking Alice's privilege with RESTRICT can be blocked by that chain.
Use CASCADE only when removing the downstream authority is the intended
administrative action. A user may still retain a privilege through a different
role or grant path, so test the effective permission after revocation.

Disabling an account is useful for preventing new authentication, but do not
assume it is equivalent to terminating every existing client session. Inspect
active sessions and close or drain them according to the operation being
performed. The public SQL surface has no general-purpose \code{KILL}
statement.

\sourcecode{src/minisql/catalog/catalog.ml}{Grant graph, role-cycle checks, and cascading revoke behavior.}
\sourcecode{src/tests/m30_audit_grants.ml}{Delegation and RESTRICT/CASCADE cases.}

\section{Use native TLS for interoperable encrypted connections}
\index{TLS}

\subsection{Certificate references and listener setup}

Windows accepts a certificate-store reference or a PFX path. The PFX password,
when needed, comes from \code{MINISQL\_TLS\_PFX\_PASSWORD}, not a
password command-line argument. Linux accepts a PEM certificate path and
private-key path separated by a vertical bar; quote the entire reference so
the shell does not interpret the bar as a pipe.

\begin{lstlisting}
# Windows: certificate/private key in a PFX file
minisqld.exe --serve-tls C:\MiniSQL-Data\db_ACTUAL_ID `
  0.0.0.0 7432 32 "pfx:C:\MiniSQL-TLS\server.pfx"

# Windows: certificate store reference
minisqld.exe --serve-tls C:\MiniSQL-Data\db_ACTUAL_ID `
  0.0.0.0 7432 32 "store:ACTUAL_CERTIFICATE_THUMBPRINT"
\end{lstlisting}

\begin{lstlisting}[language=bash]
minisqld --serve-tls /srv/minisql/db_ACTUAL_ID \
  0.0.0.0 7432 32 \
  'pem:/etc/minisql/server.pem|/etc/minisql/server-key.pem'
\end{lstlisting}

\code{0.0.0.0} listens on all IPv4 interfaces; choose a narrower address
where appropriate. The certificate must identify the name clients expect,
and the process must be able to read the private key. The fixed policy is
TLS 1.3, \code{TLS\_AES\_256\_GCM\_SHA384}, and \code{X25519}.
The configuration is validated against these names rather than offering an
arbitrary cipher-selection interface.

\begin{lstlisting}
minisql.exe --tls-shell 192.0.2.20 7432 db.example.test admin
\end{lstlisting}

The address selects the route; \code{db.example.test} is the expected
server identity. The client prompts for the database password after selecting
the TLS connection parameters. Replace documentation addresses and names with
the actual installation values.

\subsection{Trust and pinning}
\index{certificate pinning}

The native client's system-trust mode validates the X.509 chain and expected
server identity, along with certificate validity and purpose. A pin is the
SHA-256 digest of the leaf certificate. Obtain it through a trusted channel;
copying a digest from an unauthenticated connection defeats the reason for
checking it.

\begin{lstlisting}
minisql.exe --tls-pin-shell 192.0.2.20 7432 db.example.test `
  ACTUAL_64_HEX_SHA256_PIN admin
\end{lstlisting}

The native pinned path supports deliberately self-signed certificates while
retaining its native identity/validity checks. Driver options are not
identical: Python's \code{trust\_server\_certificate=True} with a required
pin, and JDBC's corresponding trust-server-certificate option, use pin-only
trust rather than the full normal chain/hostname validation path. Consult the
connector chapter before copying a native-client trust assumption into those
drivers. Normal driver TLS with an additional pin preserves the normal trust
checks.

Pinning also affects certificate renewal: a new leaf certificate has a new
digest even if it has the same hostname. Plan the trusted distribution of
updated pins. Do not silently disable verification as a routine response to
a certificate mismatch.

\sourcecode{src/minisql/platform/tls_policy.ml}{Fixed TLS policy and trust settings.}
\sourcecode{src/minisql/platform/tls_schannel.ml}{Windows native certificate handling.}
\sourcecode{src/minisql/platform/tls_openssl.ml}{Linux native TLS handling.}
\sourcecode{src/apps/minisql/main.ml}{TLS and pinned-TLS command argument order.}

\section{The MiniSQL-specific secure transport}

The \code{--serve-secure} and \code{--secure-*} families use MiniSQL's
authenticated encrypted framing. Protected frames carry an eight-byte
little-endian sequence number, ciphertext, and a 16-byte GCM tag. Associated
data binds message metadata and sequence information. This transport is not
TLS and has no X.509 certificate workflow. A normal TLS client cannot speak
to it by merely setting a TLS flag.

Use the mode explicitly when integrating an existing MiniSQL-specific secure
client. Prefer native TLS for the documented Python/JDBC TLS paths and
standards-based certificate operations. The positional authenticated loopback
mode uses the same protected frames after successful authentication, but fixes
the network address to loopback. Native TLS connections also complete database
authentication and enable this protocol protection inside the TLS channel.
Thus TLS adds a certificate-based transport boundary; it is not another name
for the MiniSQL frame format.

\sourcecode{docs/formats/secure-transport-v1.md}{Protected-frame structure.}
\sourcecode{src/minisql/server/server.ml}{Distinct authenticated, secure, and TLS listener APIs.}

\section{Enable encryption at rest}
\index{TDE}\index{encryption at rest}

MiniSQL's TDE uses a randomly generated database encryption key (DEK).
An external 32-byte key-encryption key (KEK) wraps that DEK with AES-256-GCM.
The database stores the wrapped envelope and provider identity in
\path{encryption.meta}; the external raw key file remains outside the
database tree.

Create the key directory first and protect it with OS permissions. Generate a
new key file, stop the database, and enable encryption:

\begin{lstlisting}
minisqld.exe --generate-encryption-key D:\MiniSQL-Keys\primary.key
minisqld.exe --enable-encryption D:\MiniSQL-Data\db_ACTUAL_ID `
  D:\MiniSQL-Keys\primary.key
minisql-check.exe D:\MiniSQL-Data\db_ACTUAL_ID
\end{lstlisting}

The key-generation destination must not already exist. The file contains
exactly 32 random bytes, not a human-readable password or a 64-character hex
string. The supplied path becomes the provider identifier stored in the
envelope; prefer an absolute path. Changing the process working directory
must not make a relative key reference point somewhere else.

Enabling TDE opens and closes the database, then reacquires its exclusive
process-visible lock before converting files. If a server wins that race,
the conversion fails rather than modifying an actively owned database. Each
paged file publishes its encryption feature only after the replacement is
durable, allowing the operation to resume across mixed plaintext/encrypted
files after an interruption.

Each encrypted page uses an independent nonce and authenticated data binding
it to the database identity, file identity, and page number. WAL and temporary
spill protection use separate nonce/AAD domains. This reduces the ability to
substitute a valid ciphertext from an unrelated page or file.

\warning{Coverage is artifact-specific}{TDE does not turn every byte in the
database directory into ciphertext. Text logs, schema/history side files,
audit-related artifacts, path names, deleted remnants, dumps, and filesystem
metadata require their own protection. Keep full-disk/volume encryption and
OS permissions in the deployment design. There is no current command that
disables TDE in place.}

\sourcecode{src/minisql/tools/encryption.ml}{generateKeyFile, enable, and the explicit file-conversion list.}
\sourcecode{src/minisql/security/key_provider.ml}{Raw-file provider, DEK envelopes, and path identity.}

\subsection{Rotate the external wrapping key}

Generate a different key file and rewrap the existing DEK:

\begin{lstlisting}
minisqld.exe --generate-encryption-key D:\MiniSQL-Keys\primary-next.key
minisqld.exe --rotate-encryption-key D:\MiniSQL-Data\db_ACTUAL_ID `
  D:\MiniSQL-Keys\primary-next.key
\end{lstlisting}

Rotation atomically replaces the envelope; it does not rewrite every data
page because the DEK remains the same. The implementation permits this
rewrap online. Keep the previous wrapping key until the new arrangement has
been checked and a restore drill has succeeded. Older backups may still
depend on a previous key, so key retirement must follow backup retention,
not only the current live database state.

The present provider is a raw-file provider. The tagged provider interface is
designed for additional providers, but the current executable does not
silently gain an HSM, OS keystore, or cloud KMS integration from that design.
Loss of the KEK and all usable independent backup keys means the encrypted
data cannot be recovered by resetting the SQL administrator password.

\subsection{Backups and migration need separate encryption decisions}

Use \code{backup-encrypted} with an independent backup key when a backup must
be confidential as a whole. The ordinary backup command copies the
authoritative artifacts, including a TDE envelope when present, but does not
encrypt every otherwise plaintext side file. An encrypted backup rewraps the
DEK for the backup provider and authenticates the captured files.

The page-size \code{--rewrite} migration creates a new database identity and
new files. Its copy routine does not carry over \path{encryption.meta} or
automatically enable TDE on the target. Enable and verify TDE on the new
database before treating it as having the source's at-rest protection. Also
create a new replication base after such a rewrite.

\sourcecode{src/minisql/tools/backup.ml}{runOpenEncrypted and restoreEncrypted.}
\sourcecode{src/minisql/tools/migrate.ml}{copyDatabaseState and target creation.}

\section{Audit records, ordinary logs, and SQL binlogs}
\index{audit}\index{binlog}

Audit records form a chained diagnostic history with anchors and a database
audit key. Administrative and security events use this path. Keep the audit
key, anchors, and log together in the supported backup/recovery workflow;
selectively deleting them breaks their evidentiary relationship.

Ordinary operational logs describe process activity, errors, and slow-query
summaries. The optional SQL binlog records full valid UTF-8 statements before
parsing/execution, even if the statement later fails. It is independent of the
ordinary severity threshold. A binlog write failure rejects the SQL before
execution, avoiding an unnoticed logging gap in an enabled binlog.

The three channels serve different purposes. A SQL binlog is not the physical
write-ahead log used for crash recovery, and it is not a complete logical
backup format. It can contain password literals, personal information, and
application secrets. Limit access and retention accordingly. The slow-query
record intentionally reports session, duration, row/byte counts, and success
state without including the full SQL literal payload.

\sourcecode{src/minisql/common/diagnostics.ml}{Audit records and verification.}
\sourcecode{src/minisql/common/logger.ml}{Ordinary logging and independent SQL statement capture.}

\section{An operational account design}

For a small application, create separate roles for read-only reporting,
application data changes, schema deployment, and maintenance. Give a reporting
role SELECT on only the required tables. Give an application writer INSERT,
UPDATE, and DELETE only where needed. Give schema deployment the actual
CREATE/ALTER/INDEX/DROP rights it uses. Give maintenance MAINTAIN. Reserve
ADMIN for accounts that truly need security administration, monitoring,
cancellation, or shutdown.

Current operational views require ADMIN, so a monitoring account with only
MAINTAIN cannot scrape SHOW STATUS. Treat that as an explicit privilege tradeoff
when configuring the exporter. Avoid sharing one administrator password among
all application instances simply to bypass permission setup.

Document the role graph and test it after schema changes. A newly created
table does not inherit a guessed application policy from its name. Revisit
views and procedures when their source tables or bodies change, and recreate
the replication base after security-catalog changes because the current WAL
shipping workflow is based on an archived schema/security snapshot.
